\chapter*{Záver}
\addcontentsline{toc}{chapter}{Záver}

Cieľom tejto bakalárskej práce bolo rozšíriť náš experimentálny databázový systém o ďalší dotazovací
jazyk, konkrétne o Datalog s funkčnými symbolmi a vstavanými predikátmi, a implementovať kompilátor,
ktorý bude tento jazyk prekladať do relačnej algebry. Zároveň sme chceli, aby bol výsledný
kompilátor použiteľný aj samostatne, bez priamej väzby na konkrétnu databázovú implementáciu. Tento
cieľ sa podarilo splniť.

V práci sme najprv zaviedli potrebné teoretické základy Datalogu, vysvetlili jeho sémantiku a
ukázali, aké úpravy sú potrebné na podporu funkčných symbolov a vstavaných predikátov. Dôležitou
časťou bolo aj spresnenie požiadaviek na cieľovú relačnú algebru, čo viedlo k jej predefinovaniu a
rozšíreniu o nové operátory. 

Na tomto základe sme navrhli gramatiku jazyka pre ANTLR, spôsob reprezentácie programu pomocou AST a
samotný viacstupňový algoritmus prekladu do relačnej algebry. Výsledkom je riešenie, ktoré vykonáva
nielen syntaktickú analýzu, ale aj základnú sémantickú kontrolu programu a následne konštruuje
ekvivalentný výraz relačnej algebry.

Výsledok práce je kompilátor, ktorý umožňuje používateľovi formulovať dotazy v rozšírenom Datalogu,
ktorý je prirodzenejší a stručnejší, pričom výsledný výraz relačnej algebry je generovaný
automaticky. Týmto sme dosiahli cieľ, ktorý sme si stanovili, a zároveň sme vytvorili nástroj, ktorý
môže byť ďalej rozšírený a integrovaný do rôznych databázových systémov.

Pri práci sa tiež ukázalo niekoľko obmedzení súčasného riešenia. Poradie atómov v klauzulách môže
ovplyvniť vypočítateľnosť programu, čo môže viesť k neočakávaným výsledkom. Taktiež riešenie
predpokladá existenciu vstavanej relácie $\$true$. V budúcnosti by bolo možné tieto obmedzenia
odstrániť, napríklad zavedením nového operátora na vytvorenie relácie alebo pridaním možnosti
ovplyvniť poradie operátorov v generovanom výraze relačnej algebry.

